\chapter{2012年山东卷（文）}

\section{单选题}

\begin{problem}
若复数 \(z\) 满足 \(z(2-\i)=11+7\i\)（\(\i\) 为虚数单位），则 \(z\) 为（\quad）

\choices
    {\(3+5\i\)}
    {\(3-5\i\)}
    {\(-3+5\i\)}
    {\(-3-5\i\)}
\end{problem}

\begin{answer}
A
\end{answer}

\begin{solution}
由 \(z(2-\i)=11+7\i\)，两边同乘 \(2+\i\)，得 \(5z=(11+7\i)(2+\i)=15+25\i\)，所以 \(z=3+5\i\). 故选 A.
\end{solution}

\begin{problem}
已知全集 \(U=\{0,1,2,3,4\}\)，集合 \(A=\{1,2,3\}\)，\(B=\{2,4\}\). 则 \(\left(\complement_U A\right)\cup B=\)（\quad）

\choices
    {\(\{1,2,4\}\)}
    {\(\{2,3,4\}\)}
    {\(\{0,2,4\}\)}
    {\(\{0,2,3,4\}\)}
\end{problem}

\begin{answer}
C
\end{answer}

\begin{solution}
在全集 \(U\) 中，\(A\) 的补集为 \(\complement_U A=\{0,4\}\). 因而 \(\left(\complement_U A\right)\cup B=\{0,4\}\cup\{2,4\}=\{0,2,4\}\). 故选 C.
\end{solution}

\begin{problem}
函数 \(f(x)=\dfrac{1}{\ln(x+1)}+\sqrt{4-x^2}\) 的定义域为（\quad）

\choices
    {\([-2,0)\cup(0,2]\)}
    {\((-1,0)\cup(0,2]\)}
    {\([-2,2]\)}
    {\((-1,2]\)}
\end{problem}

\begin{answer}
B
\end{answer}

\begin{solution}
根式有意义要求 \(4-x^2\geqslant0\)，即 \(-2\leqslant x\leqslant2\). 分母中的对数有意义且不为零，要求 \(x+1>0\) 且 \(\ln(x+1)\neq0\)，即 \(x>-1\) 且 \(x\neq0\). 取这些条件的交集，得到定义域为 \((-1,0)\cup(0,2]\). 故选 B.
\end{solution}

\begin{problem}
在某次测量中得到的 \(A\) 样本数据如下：82，84，84，86，86，86，88，88，88，88. 若 \(B\) 样本数据恰好是 \(A\) 样本数据每个都加 2 后所得数据，则 \(A\)，\(B\) 两样本的下列数字特征对应相同的是（\quad）

\choices
    {众数}
    {平均数}
    {中位数}
    {标准差}
\end{problem}

\begin{answer}
D
\end{answer}

\begin{solution}
把样本中每个数据都加上同一个常数 2，会使众数、平均数和中位数都增加 2. 对任意数据 \(x_i\)，新数据为 \(x_i+2\)，其平均数为 \(\overline{x}+2\)，离均差为 \((x_i+2)-(\overline{x}+2)=x_i-\overline{x}\)，所以方差和标准差均不变. 故选 D.
\end{solution}

\begin{problem}
设命题 \(p\)：函数 \(y=\sin2x\) 的最小正周期为 \(\dfrac{\pi}{2}\)；命题 \(q\)：函数 \(y=\cos x\) 的图象关于直线 \(x=\dfrac{\pi}{2}\) 对称. 则下列判断正确的是（\quad）

\choices
    {\(p\) 为真}
    {\(\neg q\) 为假}
    {\(p\wedge q\) 为假}
    {\(p\vee q\) 为真}
\end{problem}

\begin{answer}
C
\end{answer}

\begin{solution}
函数 \(y=\sin2x\) 的最小正周期为 \(\dfrac{2\pi}{2}=\pi\)，所以命题 \(p\) 为假. 对于命题 \(q\)，有 \(\cos\left(\dfrac{\pi}{2}+t\right)=-\sin t\)，而 \(\cos\left(\dfrac{\pi}{2}-t\right)=\sin t\)，二者一般不相等，因此图象不关于直线 \(x=\dfrac{\pi}{2}\) 对称，命题 \(q\) 也为假. 所以 \(p\wedge q\) 为假，选项 C 正确.
\end{solution}

\begin{problem}
设变量 \(x\)，\(y\) 满足约束条件
\(\begin{cases}
x+2y\geqslant2,\\
2x+y\leqslant4,\\
4x-y\geqslant-1,
\end{cases}\)
则目标函数 \(z=3x-y\) 的取值范围是（\quad）

\choices
    {\(\left[-\dfrac{3}{2},6\right]\)}
    {\(\left[-\dfrac{3}{2},-1\right]\)}
    {\([-1,6]\)}
    {\(\left[-6,\dfrac{3}{2}\right]\)}
\end{problem}

\begin{answer}
A
\end{answer}

\begin{solution}
可行域由三条边界直线围成. 分别联立边界方程，得到三个顶点：
\[
\begin{cases}
x+2y=2,\\
2x+y=4
\end{cases}
\Rightarrow(2,0),\qquad
\begin{cases}
x+2y=2,\\
4x-y=-1
\end{cases}
\Rightarrow(0,1),\qquad
\begin{cases}
2x+y=4,\\
4x-y=-1
\end{cases}
\Rightarrow\left(\dfrac12,3\right)
\]
目标函数为线性函数，只需比较各顶点处的值：\(z(2,0)=6\)，\(z(0,1)=-1\)，\(z\left(\dfrac12,3\right)=-\dfrac32\). 因而 \(z\in\left[-\dfrac32,6\right]\). 故选 A.
\end{solution}

\begin{problem}
执行如图所示的程序框图，如果输入 \(a=4\)，那么输出的 \(n\) 的值为（\quad）

\bitmapfigure[max width=0.55\linewidth]{img/2012/shandong_liberal/q7_fig1.png}

\choices
    {2}
    {3}
    {4}
    {5}
\end{problem}

\begin{answer}
B
\end{answer}

\begin{solution}
初始化时 \(P=0\)，\(Q=1\)，\(n=0\). 第一次判断满足 \(P\leqslant Q\)，更新为 \(P=0+4^0=1\)，\(Q=2\times1+1=3\)，\(n=1\). 第二次仍满足条件，更新为 \(P=1+4^1=5\)，\(Q=2\times3+1=7\)，\(n=2\). 第三次仍满足条件，更新为 \(P=5+4^2=21\)，\(Q=2\times7+1=15\)，\(n=3\). 此时 \(P\leqslant Q\) 不成立，输出 \(n=3\). 故选 B.
\end{solution}

\begin{problem}
函数 \(y=2\sin\left(\dfrac{\pi x}{6}-\dfrac{\pi}{3}\right)\) \(\left(0\leqslant x\leqslant9\right)\) 的最大值与最小值之和为（\quad）

\choices
    {\(2-\sqrt3\)}
    {0}
    {\(-1\)}
    {\(-1-\sqrt3\)}
\end{problem}

\begin{answer}
A
\end{answer}

\begin{solution}
令 \(t=\dfrac{\pi x}{6}-\dfrac{\pi}{3}\)，则 \(t\in\left[-\dfrac\pi3,\dfrac{7\pi}{6}\right]\). 该区间包含正弦函数取得最大值的点 \(t=\dfrac\pi2\)，所以函数最大值为 2. 正弦函数在该区间内可能取得最小值的点只有端点和 \(t=-\dfrac\pi2+2k\pi\)；后者没有落在该区间内. 两个端点处的函数值分别为 \(2\sin\left(-\dfrac\pi3\right)=-\sqrt3\) 和 \(2\sin\dfrac{7\pi}{6}=-1\)，故最小值为 \(-\sqrt3\). 因此最大值与最小值之和为 \(2-\sqrt3\). 故选 A.
\end{solution}

\begin{problem}
圆 \((x+2)^2+y^2=4\) 与圆 \((x-2)^2+(y-1)^2=9\) 的位置关系为（\quad）

\choices
    {内切}
    {相交}
    {外切}
    {相离}
\end{problem}

\begin{answer}
B
\end{answer}

\begin{solution}
两圆的圆心分别为 \(O_1(-2,0)\)、\(O_2(2,1)\)，半径分别为 \(r_1=2\)、\(r_2=3\). 圆心距为 \(d=\sqrt{(2+2)^2+(1-0)^2}=\sqrt{17}\). 因为 \(|r_2-r_1|=1<\sqrt{17}<5=r_1+r_2\)，所以两圆相交. 故选 B.
\end{solution}

\begin{problem}
函数 \(f(x)=\dfrac{\cos6x}{2^x-2^{-x}}\) 的图象大致为（\quad）

\choices
    {\choicebitmap[width=0.15\paperwidth]{img/2012/shandong_liberal/q10_optA.png}}
    {\choicebitmap[width=0.15\paperwidth]{img/2012/shandong_liberal/q10_optB.png}}
    {\choicebitmap[width=0.15\paperwidth]{img/2012/shandong_liberal/q10_optC.png}}
    {\choicebitmap[width=0.15\paperwidth]{img/2012/shandong_liberal/q10_optD.png}}
\end{problem}

\begin{answer}
D
\end{answer}

\begin{solution}
定义域为 \(\{x\in\R\mid x\ne0\}\). 又因为 \(\cos6x\) 是偶函数，而 \(2^x-2^{-x}\) 是奇函数，所以 \(f(-x)=-f(x)\)，图象关于原点对称. 当 \(x\to0^+\) 时，分母趋于 \(0^+\)，分子趋于 1，故 \(f(x)\to+\infty\)；当 \(x\to0^-\) 时，\(f(x)\to-\infty\). 当 \(|x|\to+\infty\) 时，分母绝对值趋于无穷，而分子有界，故 \(f(x)\to0\). 此外，零点满足 \(\cos6x=0\)，在两侧对称分布. 这些性质与选项 D 的图象一致，故选 D.
\end{solution}

\begin{problem}
已知双曲线 \(C_1\)：\(\dfrac{x^2}{a^2}-\dfrac{y^2}{b^2}=1\)（\(a>0\)，\(b>0\)）的离心率为 2，若抛物线 \(C_2\)：\(x^2=2py\)（\(p>0\)）的焦点到双曲线 \(C_1\) 的渐近线的距离为 2，则抛物线 \(C_2\) 的方程为（\quad）

\choices
    {\(x^2=\dfrac{8\sqrt3}{3}y\)}
    {\(x^2=\dfrac{16\sqrt3}{3}y\)}
    {\(x^2=8y\)}
    {\(x^2=16y\)}
\end{problem}

\begin{answer}
D
\end{answer}

\begin{solution}
双曲线的离心率满足 \(\dfrac{c}{a}=2\)，所以 \(c=2a\)，从而 \(b^2=c^2-a^2=3a^2\). 因而其渐近线为 \(y=\pm\sqrt3x\). 抛物线 \(x^2=2py\) 的焦点为 \(F\left(0,\dfrac p2\right)\). 取渐近线 \(\sqrt3x-y=0\)，点 \(F\) 到它的距离为 \(\dfrac{p/2}{\sqrt{3+1}}=\dfrac p4\). 由题意 \(\dfrac p4=2\)，得 \(p=8\)，所以抛物线方程为 \(x^2=16y\). 故选 D.
\end{solution}

\begin{problem}
设函数 \(f(x)=\dfrac1x\)，\(g(x)=-x^2+bx\). 若 \(y=f(x)\) 的图象与 \(y=g(x)\) 的图象有且仅有两个不同的公共点 \(A(x_1,y_1)\)，\(B(x_2,y_2)\)，则下列判断正确的是（\quad）

\choices
    {\(x_1+x_2>0,y_1+y_2>0\)}
    {\(x_1+x_2>0\)，\(y_1+y_2<0\)}
    {\(x_1+x_2<0\)，\(y_1+y_2>0\)}
    {\(x_1+x_2<0\)，\(y_1+y_2<0\)}
\end{problem}

\begin{answer}
B
\end{answer}

\begin{solution}
公共点的横坐标满足 \(\dfrac1x=-x^2+bx\)，即
\[
F(x)=x^3-bx^2+1=0
\]
该方程有两个不同的实根，故其中一个根必须是二重根. 因为 \(F'(x)=x(3x-2b)\)，且 \(F(0)=1\neq0\)，二重根只能是 \(r=\dfrac{2b}{3}\)，并且 \(F(r)=0\). 代入得 \(1-\dfrac{4b^3}{27}=0\)，所以 \(b>0\). 设另一个根为 \(s\)，由根与系数的关系，\(2r+s=b\)，从而 \(s=-\dfrac b3\). 因此两不同横坐标为 \(r\)，\(s\)，有 \(x_1+x_2=r+s=\dfrac b3>0\). 又 \(y_i=\dfrac1{x_i}\)，所以
\[
y_1+y_2=\frac{x_1+x_2}{x_1x_2}
\]
而 \(x_1x_2=rs=-\dfrac{2b^2}{9}<0\)，故 \(y_1+y_2<0\). 故选 B.
\end{solution}

\section{填空题}

\begin{problem}
如图，正方体 \(ABCD-A_1B_1C_1D_1\) 的棱长为 1，\(E\) 为线段 \(B_1C\) 上的一点，则三棱锥 \(A-DE D_1\) 的体积为\fillinblank{}.
\bitmapfigure[max width=0.48\linewidth]{img/2012/shandong_liberal/q13_fig1.png}
\end{problem}

\begin{answer}
\(\dfrac16\)
\end{answer}

\begin{solution}
取三角形 \(ADD_1\) 为底面. 由于 \(AD=DD_1=1\) 且 \(AD\perp DD_1\)，底面积为 \(\dfrac12\). 平面 \(ADD_1\) 是正方体的一个侧面，点 \(E\) 在与该平面平行且距离为 1 的侧面上，所以点 \(E\) 到平面 \(ADD_1\) 的距离为 1. 因而
\[
V_{A-DE D_1}=\frac13\times\frac12\times1=\frac16
\]
\end{solution}

\begin{problem}
如图是根据部分城市某年 6 月份的平均气温（单位：\(^{\circ}\mathrm{C}\)）数据得到的样本频率分布直方图，其中平均气温的范围是 \([20.5,26.5]\)，样本数据的分组为 \([20.5,21.5)\)，\([21.5,22.5)\)，\([22.5,23.5)\)，\([23.5,24.5)\)，\([24.5,25.5)\)，\([25.5,26.5]\).

已知样本中平均气温低于 \(22.5^{\circ}\mathrm{C}\) 的城市个数为 11，则样本中平均气温不低于 \(25.5^{\circ}\mathrm{C}\) 的城市个数为\fillinblank{}.

\bitmapfigure[max width=0.70\linewidth]{img/2012/shandong_liberal/q14_fig1.png}
\end{problem}

\begin{answer}
\(9\)
\end{answer}

\begin{solution}
前两组组距均为 1，频率分别为 \(0.10\) 和 \(0.12\)，所以平均气温低于 \(22.5^{\circ}\mathrm{C}\) 的样本频率为 \(0.10+0.12=0.22\). 设样本总数为 \(N\)，则 \(0.22N=11\)，解得 \(N=50\). 最后一组的频率为 \(0.18\times1=0.18\)，所以平均气温不低于 \(25.5^{\circ}\mathrm{C}\) 的城市个数为 \(50\times0.18=9\).
\end{solution}

\begin{problem}
若函数 \(f(x)=a^x\)（\(a>0\)，且 \(a\neq1\)）在 \([-1,2]\) 上的最大值为 4，最小值为 \(m\)，且函数 \(g(x)=(1-4m)\sqrt{x}\) 在 \((0,+\infty)\) 上是增函数，则 \(a=\)\fillinblank{}.
\end{problem}

\begin{answer}
\(\dfrac14\)
\end{answer}

\begin{solution}
当 \(a>1\) 时，\(f(x)\) 在 \([-1,2]\) 上递增，故最大值为 \(a^2=4\)，得到 \(a=2\)，最小值为 \(m=a^{-1}=\dfrac12\). 此时 \(1-4m=-1<0\)，函数 \(g(x)=-\sqrt{x}\) 在正半轴上递减，与题意矛盾.

当 \(0<a<1\) 时，\(f(x)\) 在 \([-1,2]\) 上递减，故最大值为 \(a^{-1}=4\)，得到 \(a=\dfrac14\)，最小值为 \(m=a^2=\dfrac1{16}\). 此时 \(1-4m=\dfrac34>0\)，所以 \(g(x)=\dfrac34\sqrt{x}\) 在 \((0,+\infty)\) 上递增，符合题意. 因而 \(a=\dfrac14\).
\end{solution}

\begin{problem}
如图，在平面直角坐标系 \(xOy\) 中，一单位圆的圆心初始位置在 \((0,1)\)，此时圆上一点 \(P\) 的位置在 \((0,0)\). 因圆在 \(x\) 轴上沿正向滚动，当圆滚动到圆心位于 \((2,1)\) 时，\(\overrightarrow{OP}\) 的坐标为\fillinblank{}.
\bitmapfigure[max width=0.50\linewidth]{img/2012/shandong_liberal/q16_fig1.png}
\end{problem}

\begin{answer}
\((2-\sin2,1-\cos2)\)
\end{answer}

\begin{solution}
圆心沿正向移动了 2 个单位，单位圆无滑动滚动，所以圆转过的弧长为 2，转过的圆心角为 2 弧度，方向为顺时针. 初始时从圆心指向 \(P\) 的向量为 \((0,-1)\)，顺时针旋转 2 弧度后变为 \((-\sin2,-\cos2)\). 因此点 \(P\) 的坐标为 \((2,1)+(-\sin2,-\cos2)=(2-\sin2,1-\cos2)\)，这也就是 \(\overrightarrow{OP}\) 的坐标.
\end{solution}

\section{解答题}

\begin{problem}
在 \(\triangle ABC\) 中，内角 \(A\)，\(B\)，\(C\) 所对的边分别为 \(a\)，\(b\)，\(c\)，已知 \(\sin B(\tan A+\tan C)=\tan A\tan C\).
\begin{enumerate}
\item 求证：\(a\)，\(b\)，\(c\) 成等比数列；
\item 若 \(a=1\)，\(c=2\)，求 \(\triangle ABC\) 的面积 \(S\).
\end{enumerate}
\end{problem}

\begin{solution}
\begin{enumerate}
\item 由题设中正切有意义，\(\cos A\cos C\neq0\). 又 \(A+B+C=\pi\)，所以 \(\sin(A+C)=\sin B\). 将已知等式化为
\[
\sin B\frac{\sin(A+C)}{\cos A\cos C}=\frac{\sin A\sin C}{\cos A\cos C}
\]
即 \(\sin^2B=\sin A\sin C\). 根据正弦定理，存在同一个外接圆半径 \(R\)，使 \(a=2R\sin A\)，\(b=2R\sin B\)，\(c=2R\sin C\). 因此 \(b^2=ac\)，故 \(a\)，\(b\)，\(c\) 成等比数列.
\item 由上问，\(b^2=ac=2\)，所以 \(b=\sqrt2\). 由余弦定理，\(b^2=a^2+c^2-2ac\cos B\)，从而 \(2=1+4-4\cos B\)，得 \(\cos B=\dfrac34\). 因为 \(0<B<\pi\)，所以 \(\sin B=\sqrt{1-\dfrac9{16}}=\dfrac{\sqrt7}{4}\). 因而
\[
S=\frac12ac\sin B=\frac12\times1\times2\times\frac{\sqrt7}{4}=\frac{\sqrt7}{4}
\]
\end{enumerate}
\end{solution}

\begin{problem}
袋中有五张卡片，其中红色卡片三张，标号分别为 \(1\)，\(2\)，\(3\)；蓝色卡片两张，标号分别为 \(1\)，\(2\).
\begin{enumerate}
\item 从以上五张卡片中任取两张，求这两张卡片颜色不同且标号之和小于 4 的概率；
\item 现袋中再放入一张标号为 0 的绿色卡片，从这六张卡片中任取两张，求这两张卡片颜色不同且标号之和小于 4 的概率.
\end{enumerate}
\end{problem}

\begin{solution}
\begin{enumerate}
\item 从五张卡片中任取两张的基本事件总数为 \(\mathrm C_5^2=10\). 颜色不同的取法中，按标号和小于 4 筛选，符合条件的只有红 1 与蓝 1、红 1 与蓝 2、红 2 与蓝 1，共 3 种. 所以所求概率为 \(\dfrac3{10}\).
\item 加入绿色 0 号卡片后，基本事件总数为 \(\mathrm C_6^2=15\). 原来的 3 种仍符合条件；含绿色卡片的取法中，绿色 0 号分别与红 1、红 2、红 3、蓝 1、蓝 2 配对，均颜色不同且标号和小于 4，共增加 5 种. 所以符合条件的取法共有 \(3+5=8\) 种，所求概率为 \(\dfrac8{15}\).
\end{enumerate}
\end{solution}

\begin{problem}
如图，几何体 \(E-ABCD\) 是四棱锥，\(\triangle ABD\) 为正三角形，\(CB=CD\)，\(EC\perp BD\).
\begin{enumerate}
\item 求证：\(BE=DE\)；
\item 若 \(\angle BCD=120^{\circ}\)，\(M\) 为线段 \(AE\) 的中点，求证：\(DM\parallel\) 平面 \(BEC\).
\end{enumerate}
\FigureLayoutDeclare{img/2012/shandong_liberal/q19_fig1.png}{11.0cm}{33bp 77bp 60bp 23bp}
\bitmapfigure[max width=0.44\linewidth]{img/2012/shandong_liberal/q19_fig1.png}
\end{problem}

\begin{solution}
\begin{enumerate}
\item 设 \(O\) 为 \(BD\) 的中点. 因为 \(\triangle ABD\) 为正三角形，所以 \(AB=AD\)，从而点 \(A\) 在 \(BD\) 的垂直平分面内；又因为 \(CB=CD\)，点 \(C\) 也在该垂直平分面内. 直线 \(EC\perp BD\)，且经过垂直平分面内的点 \(C\)，所以直线 \(EC\) 在该垂直平分面内，点 \(E\) 也在该平面内. 由于四棱锥的顶点 \(E\) 不在底面内，\(A\)，\(C\)，\(E\) 不共线，因而平面 \(ACE\) 就是线段 \(BD\) 的垂直平分面，所以 \(BE=DE\).
\item 取空间直角坐标系，使 \(BD\) 为 \(x\) 轴，且令 \(B=(1,0,0)\)，\(D=(-1,0,0)\). 由 \(\triangle ABD\) 为正三角形并按四边形 \(ABCD\) 的顺序取向，可设 \(A=(0,-\sqrt3,0)\). 由 \(CB=CD\) 知 \(C\) 在 \(BD\) 的垂直平分线上，设 \(C=(0,t,0)\). 由 \(\angle BCD=120^{\circ}\)，有
\[
\frac{(-1,-t)\cdot(1,-t)}{1+t^2}=\cos120^{\circ}=-\frac12
\]
从而 \(t^2=\dfrac13\). 因为四边形 \(ABCD\) 按顺序取向，点 \(A\)，\(C\) 位于对角线 \(BD\) 的两侧，所以 \(t>0\)，即 \(C=\left(0,\dfrac1{\sqrt3},0\right)\). 又因 \(EC\perp BD\)，可设 \(E=(0,u,h)\)，其中 \(u\in\R\)，\(h\neq0\). 于是
\(M=\frac{A+E}{2}=\left(0,\frac{u-\sqrt3}{2},\frac h2\right)\)，\(\overrightarrow{DM}=\left(1,\frac{u-\sqrt3}{2},\frac h2\right)\).
而 \(\overrightarrow{BC}=\left(-1,\dfrac1{\sqrt3},0\right)\)，\(\overrightarrow{CE}=\left(0,u-\dfrac1{\sqrt3},h\right)\)，所以
\[
\overrightarrow{DM}=-\overrightarrow{BC}+\frac12\overrightarrow{CE}
\]
这说明 \(\overrightarrow{DM}\) 是平面 \(BEC\) 内两个相交向量 \(\overrightarrow{BC}\)、\(\overrightarrow{CE}\) 的线性组合，故 \(DM\parallel\) 平面 \(BEC\).
\end{enumerate}
\end{solution}

\begin{problem}
已知等差数列 \(\{a_n\}\) 的前 5 项和为 105，且 \(a_{10}=2a_5\).
\begin{enumerate}
\item 求数列 \(\{a_n\}\) 的通项公式；
\item 对任意 \(m\in\N^*\)，将数列 \(\{a_n\}\) 中不大于 \(7^{2m}\) 的项的个数记为 \(b_m\). 求数列 \(\{b_m\}\) 的前 \(m\) 项和 \(S_m\).
\end{enumerate}
\end{problem}

\begin{solution}
\begin{enumerate}
\item 设等差数列的首项为 \(a_1\)，公差为 \(d\). 由前 5 项和，得 \(5(a_1+2d)=105\)，即 \(a_1+2d=21\). 又由 \(a_{10}=2a_5\)，得 \(a_1+9d=2(a_1+4d)\)，即 \(a_1=d\). 联立得 \(3a_1=21\)，所以 \(a_1=d=7\). 因而
\[
a_n=a_1+(n-1)d=7n
\]
\item 因为 \(a_n=7n\)，所以 \(a_n\leqslant7^{2m}\) 等价于 \(n\leqslant7^{2m-1}\). 其中正整数 \(n\) 的个数为 \(b_m=7^{2m-1}\). 因此
\[
S_m=\sum_{k=1}^m b_k=\sum_{k=1}^m7^{2k-1}=7\sum_{k=0}^{m-1}49^k=\frac{7(49^m-1)}{48}
\]
\end{enumerate}
\end{solution}

\begin{problem}
如图，椭圆 \(M\)：\(\dfrac{x^2}{a^2}+\dfrac{y^2}{b^2}=1\)（\(a>b>0\)）的离心率为 \(\dfrac{\sqrt3}{2}\)，直线 \(x=\pm a\) 和 \(y=\pm b\) 所围成的矩形 \(ABCD\) 的面积为 8.
\begin{enumerate}
\item 求椭圆 \(M\) 的标准方程；
\item 设直线 \(l:y=x+m\)（\(m\in\R\)）与椭圆 \(M\) 有两个不同的交点 \(P\)，\(Q\)，\(l\) 与矩形 \(ABCD\) 有两个不同的交点 \(S\)，\(T\). 求 \(\dfrac{|PQ|}{|ST|}\) 的最大值及取得最大值时 \(m\) 的值.
\end{enumerate}
\bitmapfigure[max width=0.42\linewidth]{img/2012/shandong_liberal/q21_fig1.png}
\end{problem}

\begin{solution}
\begin{enumerate}
\item 设椭圆焦半距为 \(c\). 由离心率，\(\dfrac ca=\dfrac{\sqrt3}{2}\)，所以 \(c^2=\dfrac34a^2\). 又 \(c^2=a^2-b^2\)，故 \(b^2=\dfrac14a^2\)，即 \(a=2b\). 矩形面积为 \(4ab=8\)，所以 \(ab=2\)，从而 \(b=1\)，\(a=2\). 椭圆的标准方程为
\[
\frac{x^2}{4}+y^2=1
\]
\item 直线上的两点纵坐标之差等于横坐标之差，所以同一直线上的弦长之比等于相应横坐标差之比. 将 \(y=x+m\) 代入椭圆方程，得
\[
5x^2+8mx+4m^2-4=0
\]
两交点存在且不重合的条件为 \(|m|<\sqrt5\)，两根的横坐标差为 \(\dfrac{4\sqrt{5-m^2}}5\). 矩形为 \([-2,2]\times[-1,1]\)，直线在矩形内对应的横坐标区间为 \([-2,2]\cap[-1-m,1-m]\)，故其长度为
\[
L(m)=
\begin{cases}
3-|m|,&1\leqslant|m|<\sqrt5,\\
2,&|m|\leqslant1.
\end{cases}
\]
于是
\[
R(m)=\frac{|PQ|}{|ST|}=\frac{4\sqrt{5-m^2}}{5L(m)}
\]
当 \(0\leqslant |m|\leqslant1\) 时，\(R(m)=\dfrac25\sqrt{5-m^2}\)，其最大值在 \(|m|=0\) 时取得，为 \(\dfrac{2\sqrt5}{5}\). 当 \(1\leqslant |m|<\sqrt5\) 时，令 \(t=|m|\)，则
\[
R^2(m)=\frac{16(5-t^2)}{25(3-t)^2}
\]
其导数的符号与 \(5-3t\) 相同，所以在 \(t=\dfrac53\) 时取得最大值，且
\[
R\left(\frac53\right)=\frac{2\sqrt5}{5}
\]
故最大值为 \(\dfrac{2\sqrt5}{5}\)，取得最大值时 \(m=0\) 或 \(m=\pm\dfrac53\).
\end{enumerate}
\end{solution}

\begin{problem}
已知函数 \(f(x)=\dfrac{\ln x+k}{e^x}\)（\(k\) 为常数，\(e=2.71828\cdots\) 是自然对数的底数），曲线 \(y=f(x)\) 在点 \((1,f(1))\) 处的切线与 \(x\) 轴平行.
\begin{enumerate}
\item 求 \(k\) 的值；
\item 求 \(f(x)\) 的单调区间；
\item 设 \(g(x)=xf'(x)\)，其中 \(f'(x)\) 为 \(f(x)\) 的导函数. 证明：对任意 \(x>0\)，\(g(x)<1+e^{-2}\).
\end{enumerate}
\end{problem}

\begin{solution}
\begin{enumerate}
\item 函数定义域为 \((0,+\infty)\)，且
\[
f'(x)=e^{-x}\left(\frac1x-\ln x-k\right)
\]
切线与 \(x\) 轴平行意味着 \(f'(1)=0\)，所以 \(1-k=0\)，得 \(k=1\).
\item 当 \(k=1\) 时，令 \(h(x)=\dfrac1x-\ln x-1\). 有 \(h(1)=0\)，且 \(h'(x)=-\dfrac1{x^2}-\dfrac1x<0\). 因而 \(h(x)>0\) 在 \((0,1)\) 上成立，\(h(x)<0\) 在 \((1,+\infty)\) 上成立. 由于 \(e^{-x}>0\)，所以 \(f(x)\) 在 \((0,1)\) 上单调递增，在 \((1,+\infty)\) 上单调递减.
\item 由 \(k=1\)，有
\[
g(x)=xf'(x)=e^{-x}(1-x\ln x-x)
\]
当 \(x\geqslant1\) 时，由上问知 \(f'(x)\leqslant0\)，所以 \(g(x)\leqslant0<1+e^{-2}\). 当 \(0<x<1\) 时，令 \(F(x)=1-x\ln x-x\). 则 \(F'(x)=-(\ln x+2)\)，所以 \(F\) 在 \((0,e^{-2})\) 上递增，在 \((e^{-2},1)\) 上递减，且 \(F(e^{-2})=1+e^{-2}\). 因而 \(F(x)\leqslant1+e^{-2}\). 又 \(\lim_{x\to0^+}F(x)=1\)，且 \(F(1)=0\)，结合上述单调性知 \(F(x)>0\)；又因 \(e^{-x}<1\)，故 \(g(x)=e^{-x}F(x)<F(x)\leqslant1+e^{-2}\). 综上，对任意 \(x>0\)，都有 \(g(x)<1+e^{-2}\).
\end{enumerate}
\end{solution}
